\chapter{Worked examples: further symmetric-key results}

Six smaller modules complete the symmetric-key part of Rosulek, \emph{The Joy of
Cryptography}: the insecurity of deterministic encryption (Chapter~7), a PRG from
a PRF (Chapter~6), a MAC from a PRF (Chapter~10), universal hashing and the
Carter--Wegman MAC (Chapter~12), a standalone IND-CPA game, and a hybrid argument
over the deep embedding. Several of them are stated over $\Bool$ so that every
probability is computed exactly; the text below records where a statement is
narrower than the textbook claim.

\section{Deterministic encryption is not CPA-secure}

\begin{theorem}[Deterministic XOR encryption is broken]
  \label{thm:detcpa-insecure}
  \uses{def:advantage}
  \lean{CatCryptCore.Examples.DetCPA.same_msg_prTrue, CatCryptCore.Examples.DetCPA.diff_msg_prTrue, CatCryptCore.Examples.DetCPA.det_enc_xor_insecure}
  \leanok
  For the deterministic scheme $\mathsf{enc}(k, m) = k \oplus m$, the game that
  encrypts $m$ twice and the game that encrypts $m$ and $\lnot m$ are
  distinguished with advantage exactly $1$ by the adversary that compares the two
  ciphertexts (Rosulek, §7.1).
\end{theorem}

\section{A standalone IND-CPA game}

\begin{definition}[IND-CPA games]
  \label{def:indcpa-standalone}
  \uses{def:spcomp, def:advantage}
  \lean{CatCryptCore.Examples.INDCPA.indCPA_real, CatCryptCore.Examples.INDCPA.INDCPA_Advantage, CatCryptCore.Examples.INDCPA.PerfectINDCPA}
  \leanok
  The left and right challenge games of a symmetric scheme, their advantage, and
  perfect IND-CPA security as a relational judgement between them.
\end{definition}

\begin{theorem}[Perfect IND-CPA gives advantage zero]
  \label{thm:indcpa-standalone-zero}
  \uses{def:indcpa-standalone, thm:advantage-zero-of-rhoare}
  \lean{CatCryptCore.Examples.INDCPA.perfect_indcpa_zero_advantage}
  \leanok
  A perfectly IND-CPA scheme has advantage $0$ against every adversary.
\end{theorem}

\section{A PRG from a PRF}

Rosulek, Claim~6.3: $G(s) = (F(s, 0), F(s, 1))$. The module works over
$\mathsf{Word} = \Bool$ with two counter values.

\begin{definition}[Decomposed hybrids]
  \label{def:prfprg-decomp}
  \uses{def:spcomp}
  \lean{CatCrypt.Examples.PRFPRG.gen_hyb_decomp, CatCrypt.Examples.PRFPRG.prf_advantage}
  \leanok
  Hybrid games in which the two PRF evaluations use independent seeds; hybrid
  $i$ replaces the first $i$ outputs by uniform words.
\end{definition}

\begin{theorem}[Hybrid bound]
  \label{thm:prfprg-bound}
  \uses{def:prfprg-decomp, thm:advantageA-triangle}
  \lean{CatCrypt.Examples.PRFPRG.prg_decomp_security_bound}
  \leanok
  The distance between the first and last decomposed hybrid is at most the sum of
  two PRF advantages of explicit reduction adversaries. The bound concerns the
  decomposed games, in which the two evaluations use independent seeds; it is not
  a bound on the counter-mode generator with a single seed.
\end{theorem}

\begin{theorem}[The XOR instantiation]
  \label{thm:prfprg-xor}
  \uses{def:prfprg-decomp}
  \lean{CatCrypt.Examples.PRFPRG.xorPRF_zero_advantage, CatCrypt.Examples.PRFPRG.xorPRF_hyb_1_to_2_zero_advantage}
  \leanok
  For the XOR function the single-query PRF advantage is $0$, and so is the
  advantage between the second and third hybrids of the single-seed construction.
\end{theorem}

\section{A MAC from a PRF}

Rosulek, Claim~10.4: $\mathsf{GetTag}(m) = F_k(m)$. The module proves the
individual game hops over $\Bool$ and composes them for a forgery on a message
other than the one queried.

\begin{lemma}[Game hops]
  \label{lem:prfmac-hops}
  \uses{lem:liftr-uniform-bij}
  \lean{CatCrypt.Examples.PRFMAC.prf_uniform, CatCrypt.Examples.PRFMAC.mac_same_message_coupling, CatCrypt.Examples.PRFMAC.forgery_diff_message_uniform}
  \leanok
  The PRF output on a uniform key is uniform; on the queried message the real and
  ideal tags couple; a forgery on another message verifies with a uniform bit.
\end{lemma}

\begin{theorem}[Forgery bound]
  \label{thm:prfmac-bound}
  \uses{lem:prfmac-hops, thm:advantageA-triangle}
  \lean{CatCrypt.Examples.PRFMAC.security_based_on_prf}
  \leanok
  For a forgery on a message different from the queried one, the distance between
  the real and ideal MAC games is at most two PRF advantages plus a statistical
  gap of $1/2$, the guessing probability of a one-bit tag.
\end{theorem}

\begin{theorem}[No perfect coupling on an unqueried message]
  \label{thm:prfmac-no-coupling}
  \uses{lem:prfmac-hops}
  \lean{CatCrypt.Examples.PRFMAC.not_rHoare_macReal_macIdeal_of_ne}
  \leanok
  For a forgery on a message different from the queried one, the real and ideal
  MAC games admit no equality coupling: the identity distinguisher has advantage
  $1/2$, so the statistical gap in the forgery bound cannot be removed.
\end{theorem}

\section{Universal hashing and the Carter--Wegman MAC}

\begin{definition}[Universal hash family]
  \label{def:uhf}
  \uses{def:spcomp}
  \lean{CatCryptCore.Examples.UniversalHash.UHFamily, CatCryptCore.Examples.UniversalHash.IsUniversal, CatCryptCore.Examples.UniversalHash.boolPairwiseHash}
  \leanok
  A keyed hash family with a collision bound $\varepsilon$; the pairwise
  independent family $H_{k_1,k_2}(m) = (k_1 \wedge m) \oplus k_2$ over $\Bool$.
\end{definition}

\begin{theorem}[Universality and one-time MAC security]
  \label{thm:uhf-mac}
  \uses{def:uhf}
  \lean{CatCryptCore.Examples.UniversalHash.boolPairwiseHash_is_universal, CatCryptCore.Examples.UniversalHash.boolHash_forge_advantage, CatCryptCore.Examples.UniversalHash.boolPairwiseHash_mac_security}
  \leanok
  The $\Bool$ family is $1/2$-universal, and the one-time forgery advantage on a
  message other than the queried one is at most $1/2$.
\end{theorem}

\section{A hybrid argument for XOR encryption}

\begin{theorem}[Two- and three-instance hybrids]
  \label{thm:deephybrid}
  \uses{thm:advantageA-triangle}
  \lean{CatCrypt.Examples.DeepHybrid.single_oracle_zero_advantage, CatCrypt.Examples.DeepHybrid.two_instance_advantage_bound, CatCrypt.Examples.DeepHybrid.three_instance_triangle, CatCrypt.Examples.DeepHybrid.hybrid_step_zero_advantage}
  \leanok
  For XOR encryption with independent keys, a single instance has advantage $0$;
  the two-instance game is bounded by its two hybrid steps, the three-instance
  game by the sum of its three steps, and each step of the three-instance ladder
  has advantage $0$.
\end{theorem}
